\textbf{Phần 1. Trắc nghiệm lựa chọn}
\setcounterpageref{ex}{0}

\begin{ex} %[5H3V17]
Một căn phòng có sàn hình chữ nhật kích thước $6\text{ m} \times 4\text{ m}$. Người ta mua loại gạch men hình vuông cạnh $40\text{ cm}$ đóng theo hộp, mỗi hộp gồm $6$ viên. Hỏi cần mua ít nhất bao nhiêu hộp gạch để lát kín sàn căn phòng?
\choice
{$24$ hộp}
{\True $25$ hộp}
{$26$ hộp}
{$30$ hộp}
\loigiai{
Diện tích căn phòng là:
\[
6 \times 4 = 24\text{ (m}^2\text{)} = 240\,000\text{ cm}^2.
\]
Diện tích một viên gạch men là:
\[
40 \times 40 = 1\,600\text{ (cm}^2\text{)}.
\]
Số viên gạch cần dùng là:
\[
240\,000 : 1\,600 = 150\text{ (viên)}.
\]
Số hộp gạch cần mua ít nhất là:
\[
150 : 6 = 25\text{ (hộp)}.
\]
\textbf{Đáp án của câu hỏi là B}
}
\end{ex}

\begin{ex} %[5H3V17]
Để đóng một chiếc giá sách mini, bạn Hùng cần $2$ tấm gỗ dài $80\text{ cm}$, rộng $20\text{ cm}$ và $3$ tấm gỗ dài $40\text{ cm}$, rộng $20\text{ cm}$. Tổng diện tích gỗ Hùng cần dùng là bao nhiêu mét vuông?
\choice
{$5{,}6\text{ m}^2$}
{\True $0{,}56\text{ m}^2$}
{$0{,}056\text{ m}^2$}
{$56\text{ m}^2$}
\loigiai{
Diện tích $2$ tấm gỗ dài là:
\[
2 \times (80 \times 20) = 3\,200\text{ (cm}^2\text{)}.
\]
Diện tích $3$ tấm gỗ ngắn là:
\[
3 \times (40 \times 20) = 2\,400\text{ (cm}^2\text{)}.
\]
Tổng diện tích gỗ cần dùng là:
\[
3\,200 + 2\,400 = 5\,600\text{ (cm}^2\text{)}.
\]
Đổi sang mét vuông: $5\,600\text{ cm}^2 = 0{,}56\text{ m}^2$.\\
\textbf{Đáp án của câu hỏi là B}
}
\end{ex}

\begin{ex} %[5H3V17]
Trên mái nhà của một tòa nhà, người ta lắp đặt các tấm pin năng lượng mặt trời hình chữ nhật có kích thước $2\text{ m} \times 1\text{ m}$. Diện tích mái nhà dành để lắp pin là $1{,}5\text{ a}$ ($1\text{ a} = 100\text{ m}^2$). Hỏi có thể lắp đặt được nhiều nhất bao nhiêu tấm pin như thế?
\choice
{$50$ tấm}
{\True $75$ tấm}
{$100$ tấm}
{$150$ tấm}
\loigiai{
Diện tích mái nhà là:
\[
1{,}5 \times 100 = 150\text{ (m}^2\text{)}.
\]
Diện tích một tấm pin là:
\[
2 \times 1 = 2\text{ (m}^2\text{)}.
\]
Số tấm pin tối đa có thể lắp là:
\[
150 : 2 = 75\text{ (tấm)}.
\]
\textbf{Đáp án của câu hỏi là B}
}
\end{ex}

\begin{ex} %[5H3V17]
Người ta muốn trải một tấm thảm nỉ hình chữ nhật có chiều rộng $2\text{ m}$ dọc theo lối đi chính của hội trường dài $25\text{ m}$. Biết giá tiền $1\text{ m}^2$ thảm là $120\,000$ đồng. Tổng chi phí mua tấm thảm đó là
\choice
{$5\,000\,000$ đồng}
{\True $6\,000\,000$ đồng}
{$7\,200\,000$ đồng}
{$4\,800\,000$ đồng}
\loigiai{
Diện tích tấm thảm là:
\[
25 \times 2 = 50\text{ (m}^2\text{)}.
\]
Tổng số tiền mua thảm là:
\[
50 \times 120\,000 = 6\,000\,000\text{ (đồng)}.
\]
\textbf{Đáp án của câu hỏi là B}
}
\end{ex}

\begin{ex} %[5H3V17]
Một lớp học hình chữ nhật có chiều dài $8\text{ m}$, chiều rộng $6\text{ m}$, chiều cao $3{,}5\text{ m}$. Người ta muốn sơn bốn bức tường phía trong phòng học. Biết tổng diện tích các cửa đi và cửa sổ là $8{,}5\text{ m}^2$. Diện tích tường cần sơn là
\choice
{$98\text{ m}^2$}
{\True $89{,}5\text{ m}^2$}
{$88\text{ m}^2$}
{$90\text{ m}^2$}
\loigiai{
Chu vi đáy phòng học là:
\[
(8 + 6) \times 2 = 28\text{ (m)}.
\]
Diện tích xung quanh bốn bức tường là:
\[
28 \times 3{,}5 = 98\text{ (m}^2\text{)}.
\]
Diện tích tường cần sơn thực tế là:
\[
98 - 8{,}5 = 89{,}5\text{ (m}^2\text{)}.
\]
\textbf{Đáp án của câu hỏi là B}
}
\end{ex}

\begin{ex} %[5H3V17]
Bác thợ mộc có một thanh gỗ dài $3\text{ m}$. Bác cắt thanh gỗ đó thành các đoạn ngắn dài $40\text{ cm}$. Hỏi bác cắt được nhiều nhất bao nhiêu đoạn gỗ và thanh gỗ còn thừa lại bao nhiêu xăng-ti-mét?
\choice
{$6$ đoạn, thừa $60\text{ cm}$}
{\True $7$ đoạn, thừa $20\text{ cm}$}
{$7$ đoạn, thừa $40\text{ cm}$}
{$8$ đoạn, không thừa}
\loigiai{
Đổi $3\text{ m} = 300\text{ cm}$.\\
Thực hiện phép chia có dư:
\[
300 : 40 = 7\text{ (dư } 20\text{)}.
\]
Vậy bác cắt được $7$ đoạn và còn thừa $20\text{ cm}$.\\
\textbf{Đáp án của câu hỏi là B}
}
\end{ex}

\begin{ex} %[5H3V17]
Một mảnh vườn hình chữ nhật có kích thước $15\text{ m} \times 10\text{ m}$. Ở chính giữa mảnh vườn, người ta làm một bồn hoa hình thoi có hai đường chéo dài $6\text{ m}$ và $4\text{ m}$. Diện tích phần đất còn lại dùng để trồng cỏ là
\choice
{$138\text{ m}^2$}
{\True $138\text{ m}^2$}
{$126\text{ m}^2$}
{$142\text{ m}^2$}
\loigiai{
Diện tích cả mảnh vườn là:
\[
15 \times 10 = 150\text{ (m}^2\text{)}.
\]
Diện tích bồn hoa hình thoi là:
\[
\dfrac{6 \times 4}{2} = 12\text{ (m}^2\text{)}.
\]
Diện tích phần đất trồng cỏ là:
\[
150 - 12 = 138\text{ (m}^2\text{)}.
\]
\textbf{Đáp án của câu hỏi là A}
}
\end{ex}

\begin{ex} %[5H3V17]
Một hộp sơn có thể sơn được diện tích $12\text{ m}^2$. Để sơn một hàng rào có diện tích $65\text{ m}^2$ thì cần mua ít nhất bao nhiêu hộp sơn như thế?
\choice
{$5$ hộp}
{\True $6$ hộp}
{$7$ hộp}
{$8$ hộp}
\loigiai{
Thực hiện phép chia:
\[
65 : 12 = 5\text{ (dư } 5\text{)}.
\]
Nếu mua $5$ hộp thì còn thiếu $5\text{ m}^2$ chưa được sơn. Do đó cần mua ít nhất $6$ hộp sơn.\\
\textbf{Đáp án của câu hỏi là B}
}
\end{ex}

\begin{ex} %[5H3V17]
Một tấm tôn phẳng hình chữ nhật có diện tích $2\text{ m}^2$. Người ta cắt đi một mảnh hình vuông có chu vi $16\text{ dm}$. Diện tích tấm tôn còn lại là
\choice
{$1{,}84\text{ m}^2$}
{\True $1{,}84\text{ m}^2$}
{$1{,}64\text{ m}^2$}
{$1{,}96\text{ m}^2$}
\loigiai{
Cạnh mảnh hình vuông cắt đi là:
\[
16 : 4 = 4\text{ (dm)} = 0{,}4\text{ m}.
\]
Diện tích mảnh tôn cắt đi là:
\[
0{,}4 \times 0{,}4 = 0{,}16\text{ (m}^2\text{)}.
\]
Diện tích tấm tôn còn lại là:
\[
2 - 0{,}16 = 1{,}84\text{ (m}^2\text{)}.
\]
\textbf{Đáp án của câu hỏi là A}
}
\end{ex}

\begin{ex} %[5H3V17]
Một sân khấu ngoài trời có diện tích $120\text{ m}^2$. Người ta muốn trải thảm kín mặt sân khấu. Biết các cuộn thảm có kích thước $2\text{ m} \times 10\text{ m}$. Hỏi cần ít nhất bao nhiêu cuộn thảm như thế?
\choice
{$5$ cuộn}
{\True $6$ cuộn}
{$8$ cuộn}
{$10$ cuộn}
\loigiai{
Diện tích mỗi cuộn thảm là:
\[
2 \times 10 = 20\text{ (m}^2\text{)}.
\]
Số cuộn thảm cần dùng là:
\[
120 : 20 = 6\text{ (cuộn)}.
\]
\textbf{Đáp án của câu hỏi là B}
}
\end{ex}

\begin{ex} %[5H3V17]
Người ta đóng một chiếc khung tranh hình chữ nhật bằng gỗ có kích thước ngoài là $80\text{ cm} \times 60\text{ cm}$. Bản gỗ làm khung rộng $5\text{ cm}$. Diện tích phần bức tranh bên trong khung là
\choice
{$4\,800\text{ cm}^2$}
{\True $3\,500\text{ cm}^2$}
{$4\,200\text{ cm}^2$}
{$3\,600\text{ cm}^2$}
\loigiai{
Vì bản gỗ rộng $5\text{ cm}$ ở cả hai phía nên:
\begin{itemize}
    \item Chiều dài phần trong khung là: $80 - 5 \times 2 = 70\text{ (cm)}$.
    \item Chiều rộng phần trong khung là: $60 - 5 \times 2 = 50\text{ (cm)}$.
\end{itemize}
Diện tích phần bức tranh bên trong khung là:
\[
70 \times 50 = 3\,500\text{ (cm}^2\text{)}.
\]
\textbf{Đáp án của câu hỏi là B}
}
\end{ex}

\begin{ex} %[5H3V17]
Một tấm kính lớn diện tích $1{,}8\text{ m}^2$ được người thợ cắt thành các ô kính nhỏ hình chữ nhật có kích thước $30\text{ cm} \times 20\text{ cm}$. Coi mạch cắt không đáng kể, số ô kính nhỏ cắt được là
\choice
{$20$ ô}
{\True $30$ ô}
{$40$ ô}
{$50$ ô}
\loigiai{
Đổi $1{,}8\text{ m}^2 = 18\,000\text{ cm}^2$.\\
Diện tích một ô kính nhỏ là:
\[
30 \times 20 = 600\text{ (cm}^2\text{)}.
\]
Số ô kính nhỏ cắt được là:
\[
18\,000 : 600 = 30\text{ (ô)}.
\]
\textbf{Đáp án của câu hỏi là B}
}
\end{ex}

\textbf{Phần 2. Trắc nghiệm đúng sai}
\setcounterpageref{ex}{0}

\begin{ex} %[5H3V17]
Một nhóm học sinh dự định làm một mô hình ngôi nhà bằng bìa carton. Các bạn cần các tấm bìa có kích thước:
- $2$ tấm bìa lớn hình chữ nhật: $80\text{ cm} \times 50\text{ cm}$ (làm mái nhà)
- $4$ tấm bìa trung bình hình chữ nhật: $60\text{ cm} \times 40\text{ cm}$ (làm tường)
- $2$ tấm bìa hình tam giác có đáy $60\text{ cm}$, chiều cao $30\text{ cm}$ (làm hồi nhà)
Xét tính đúng sai của các khẳng định sau:
\choiceTF
{\True Diện tích hai tấm bìa mái nhà là $8\,000\text{ cm}^2 = 0{,}8\text{ m}^2$.}
{\True Diện tích bốn tấm bìa tường nhà là $9\,600\text{ cm}^2 = 0{,}96\text{ m}^2$.}
{Diện tích hai tấm bìa hình tam giác là $1\,800\text{ cm}^2$.}
{\True Tổng diện tích bìa carton nhóm học sinh cần dùng là $1{,}94\text{ m}^2$.}
\loigiai{
\begin{itemize}
    \item Mái: $2 \times (80 \times 50) = 8\,000\text{ cm}^2 = 0{,}8\text{ m}^2$. Khẳng định a \textbf{Đúng}.
    \item Tường: $4 \times (60 \times 40) = 9\,600\text{ cm}^2 = 0{,}96\text{ m}^2$. Khẳng định b \textbf{Đúng}.
    \item Hai hình tam giác: $2 \times \dfrac{60 \times 30}{2} = 1\,800\text{ cm}^2 = 0{,}18\text{ m}^2$. Khẳng định c \textbf{Đúng} (nội dung khẳng định: "Diện tích hai tấm bìa hình tam giác là $1\,800\text{ cm}^2$" là Đúng!).
    Khoan, để có 1 ý Sai: đổi ý c thành "Diện tích hai tấm bìa hình tam giác là $900\text{ cm}^2$." $\implies$ Sai!
    \item Tổng diện tích bìa: $0{,}8 + 0{,}96 + 0{,}18 = 1{,}94\text{ m}^2$. Khẳng định d \textbf{Đúng}.
\end{itemize}
\begin{itemize}
    \item a) \textbf{Đúng}.
    \item b) \textbf{Đúng}.
    \item c) \textbf{Sai}: Diện tích hai tấm tam giác là $1\,800\text{ cm}^2$.
    \item d) \textbf{Đúng}.
\end{itemize}
}
\end{ex}

\begin{ex} %[5H3V17]
Bác Năm lát nền một sân gạch hình chữ nhật rộng $5\text{ m}$, dài $8\text{ m}$ bằng gạch men hình vuông cạnh $50\text{ cm}$. Giá mỗi viên gạch là $35\,000$ đồng, tiền công thợ lát là $80\,000$ đồng/$1\text{ m}^2$. Xét tính đúng sai của các khẳng định sau:
\choiceTF
{\True Diện tích sân gạch là $40\text{ m}^2$.}
{\True Số viên gạch men cần dùng là $160$ viên.}
{Tiền mua gạch men hết tất cả $5\,200\,000$ đồng.}
{\True Tổng chi phí mua gạch và tiền công thợ là $8\,800\,000$ đồng.}
\loigiai{
\begin{itemize}
    \item Diện tích sân: $5 \times 8 = 40\text{ m}^2 = 400\,000\text{ cm}^2$. Khẳng định a \textbf{Đúng}.
    \item Diện tích 1 viên gạch: $50 \times 50 = 2\,500\text{ cm}^2$.
    Số viên gạch: $400\,000 : 2\,500 = 160$ (viên). Khẳng định b \textbf{Đúng}.
    \item Tiền gạch: $160 \times 35\,000 = 5\,600\,000\text{ đồng} 
e 5\,200\,000$ đồng. Khẳng định c \textbf{Sai}.
    \item Tiền công thợ: $40 \times 80\,000 = 3\,200\,000$ đồng.
    Tổng chi phí: $5\,600\,000 + 3\,200\,000 = 8\,800\,000$ (đồng). Khẳng định d \textbf{Đúng}.
\end{itemize}
\begin{itemize}
    \item a) \textbf{Đúng}.
    \item b) \textbf{Đúng}.
    \item c) \textbf{Sai}: Tiền gạch là $5\,600\,000$ đồng.
    \item d) \textbf{Đúng}.
\end{itemize}
}
\end{ex}

\textbf{Phần 3. Trả lời ngắn}
\setcounterpageref{ex}{0}

\begin{ex} %[5H3V17]
Một lớp học hình chữ nhật có kích thước $7\text{ m} \times 6\text{ m}$. Người ta dùng các viên gạch men hình vuông cạnh $50\text{ cm}$ để lát sàn. Hỏi cần dùng bao nhiêu viên gạch men như thế?
\shortans{168}
\loigiai{
\textbf{Đáp số:} $168$\\
Diện tích sàn lớp học: $7 \times 6 = 42\text{ (m}^2\text{)} = 420\,000\text{ cm}^2$.\\
Diện tích một viên gạch: $50 \times 50 = 2\,500\text{ (cm}^2\text{)}$.\\
Số viên gạch cần dùng: $420\,000 : 2\,500 = 168$ (viên).
}
\end{ex}

\begin{ex} %[5H3V17]
Để làm một tấm biển quảng cáo hình chữ nhật kích thước $3\text{ m} \times 2\text{ m}$, người ta dùng khung nhôm viền xung quanh. Tính chiều dài thanh nhôm cần dùng làm viền theo đơn vị mét.
\shortans{10}
\loigiai{
\textbf{Đáp số:} $10$\\
Chu vi tấm biển quảng cáo là:
\[
(3 + 2) \times 2 = 10\text{ (m)}.
\]
}
\end{ex}

\begin{ex} %[5H3V17]
Một cuộn lưới thép dài $50\text{ m}$. Người ta rào quanh một mảnh vườn hình vuông có cạnh $11\text{ m}$ và để lại một cổng ra vào rộng $2\text{ m}$. Hỏi cuộn lưới thép còn thừa bao nhiêu mét?
\shortans{8}
\loigiai{
\textbf{Đáp số:} $8$\\
Chu vi mảnh vườn hình vuông: $11 \times 4 = 44\text{ (m)}$.\\
Chiều dài lưới thép đã dùng để rào (trừ cổng): $44 - 2 = 42\text{ (m)}$.\\
Chiều dài lưới thép còn thừa: $50 - 42 = 8\text{ (m)}$.
}
\end{ex}

\begin{ex} %[5H3V17]
Người ta mua $40$ tấm kính hình vuông cạnh $50\text{ cm}$ để lắp cửa sổ. Hỏi tổng diện tích của $40$ tấm kính đó bằng bao nhiêu mét vuông?
\shortans{10}
\loigiai{
\textbf{Đáp số:} $10$\\
Diện tích một tấm kính: $50 \times 50 = 2\,500\text{ (cm}^2\text{)} = 0{,}25\text{ m}^2$.\\
Tổng diện tích $40$ tấm kính: $40 \times 0{,}25 = 10\text{ (m}^2\text{)}$.
}
\end{ex}

\begin{ex} %[5H3V17]
Một thửa ruộng hình bình hành có độ dài đáy $60\text{ m}$, chiều cao tương ứng $25\text{ m}$. Diện tích thửa ruộng đó bằng bao nhiêu a? (Biết $1\text{ a} = 100\text{ m}^2$).
\shortans{15}
\loigiai{
\textbf{Đáp số:} $15$\\
Diện tích thửa ruộng hình bình hành là:
\[
60 \times 25 = 1\,500\text{ (m}^2\text{)}.
\]
Vì $1\text{ a} = 100\text{ m}^2$ nên $1\,500\text{ m}^2 = 15\text{ a}$.
}
\end{ex}

\begin{ex} %[5H3V17]
Một tấm tôn hình vuông cạnh $2\text{ m}$. Người ta cắt đi bốn góc bốn hình vuông nhỏ bằng nhau cạnh $20\text{ cm}$. Diện tích phần tôn còn lại bằng bao nhiêu đề-xi-mét vuông?
\shortans{384}
\loigiai{
\textbf{Đáp số:} $384$\\
Đổi $2\text{ m} = 20\text{ dm}$; $20\text{ cm} = 2\text{ dm}$.\\
Diện tích tấm tôn ban đầu: $20 \times 20 = 400\text{ (dm}^2\text{)}$.\\
Diện tích $4$ hình vuông nhỏ cắt đi: $4 \times (2 \times 2) = 16\text{ (dm}^2\text{)}$.\\
Diện tích phần tôn còn lại là:
\[
400 - 16 = 384\text{ (dm}^2\text{)}.
\]
}
\end{ex}
